\part{The shell}

\chapter{The read-eval-print loop}

\begin{goals}
  \item How a command interpreter is structured
  \item Follow a command from keystroke to effect
  \item See how every subsystem in the course is used at once
\end{goals}

\section{The whole program}

\filetag{lytlnybl/user/programs/shell.c}
\begin{lstlisting}[style=c]
void start(void) {
  for (;;) {
    printf("> ");                       /* Print the prompt */

    char input[50];
    read(0, input, 50);                 /* Read */
    input[strlen(input) - 1] = '\0';

    size_t word_count;
    char** words = tokenize_line(input, &word_count);   /* Eval */

    if (!execute_command(words, word_count)) {
      printf("command failed\n");
    }

    memset(input, 0, strlen(input));    /* Loop */
    word_count = 0;
    free(words);
  }
}
\end{lstlisting}

That is a \term{REPL}: Read, Eval, Print, Loop. Every interactive interpreter ever written --- bash,
Python, a database console --- has this shape.

\begin{concept}{Everything you built, in eleven lines}
Look at what those lines depend on:
\begin{itemize}
  \item \cod{printf} $\rightarrow$ \cod{write} $\rightarrow$ \cod{int 0x80} $\rightarrow$ the IDT
        $\rightarrow$ the VGA driver.
  \item \cod{read} $\rightarrow$ the process blocks $\rightarrow$ the scheduler switches away
        $\rightarrow$ a key arrives $\rightarrow$ IRQ1 $\rightarrow$ the keyboard driver writes into
        another address space $\rightarrow$ the process is woken.
  \item \cod{malloc} inside \cod{tokenize\_line} $\rightarrow$ \cod{sbrk} $\rightarrow$ the kernel
        maps a page $\rightarrow$ \cod{alloc\_frame}.
  \item \cod{execute\_command} $\rightarrow$ filesystem calls $\rightarrow$ path resolution
        $\rightarrow$ inodes $\rightarrow$ the ATA driver $\rightarrow$ I/O ports.
\end{itemize}
Every part of this course is exercised by one iteration of this loop.
\end{concept}

\section{Tokenising}

\filetag{lytlnybl/user/programs/shell.c}
\begin{lstlisting}[style=c]
char** tokenize_line(const char* line, size_t *count_out) {
  const char* p = line;
  char **words;
  size_t words_i = 0;
  size_t capacity = 8;

  words = calloc(capacity + 1, sizeof(char *));

  while (*p != '\0') {
    (*count_out)++;
    const char* start;

    while (isspace((unsigned char) *p)) p++;
    if (*p == '\0') break;

    start = p;
    while (*p != '\0' && !isspace((unsigned char) *p)) p++;

    size_t len = p - start;
    char *word = malloc(len + 1);
    if (word == NULL) return NULL;

    memcpy(word, start, len);
    word[len] = '\0';
    words[words_i++] = word;
  }
  return words;
}
\end{lstlisting}

The algorithm is right: skip whitespace, mark the start, advance to the next whitespace, copy the
slice. But the bookkeeping around it has three genuine defects.

\begin{danger}{Three bugs in the word counter}
\textbf{1. \cod{word\_count} is never initialised.} In \cod{start()} it is declared as
\cod{size\_t word\_count;} and passed by address. \cod{tokenize\_line} then does
\cod{(*count\_out)++} on whatever garbage was on the stack. Since \cod{execute\_command} tests
\cod{word\_count != 2}, the shell's behaviour depends on an indeterminate value.

It works in practice only because the stack slot happens to be zero on the first iteration and the
variable is reset to 0 at the end of each loop --- which is itself an accident, because the reset
happens \emph{after} the value has already been used.

\textbf{2. The increment is in the wrong place.} \cod{(*count\_out)++} runs at the top of the loop,
before checking whether a token actually exists. A trailing space makes the loop run one extra time
and count a word that was never produced.

\textbf{3. \cod{capacity} is never enforced.} The array holds 9 pointers. A line with 12 words
writes past the end of a heap allocation, corrupting the next block's header --- which, per Chapter
29, corrupts the allocator itself.
\end{danger}

\begin{warn}{And a leak}
\cod{start()} calls \cod{free(words)}, which releases the array of pointers --- but every string
inside it was separately \cod{malloc}ed and is never freed. Each command leaks a few dozen bytes.
With a 4 KB heap that grows by \cod{sbrk}, the shell slowly consumes memory for as long as it runs.

The correct cleanup is a loop:
\begin{lstlisting}[style=c]
for (size_t i = 0; i < word_count; i++) free(words[i]);
free(words);
\end{lstlisting}
\end{warn}

\begin{concept}{Why none of this is surprising}
These are not sloppy mistakes; they are the normal state of code that has not yet been stress
tested. The shell is the last thing built in the project, it works for the inputs it was tried with,
and the bugs only appear at the edges: an empty line, a trailing space, twelve words, a long
session.

The lesson worth taking is the one about \emph{compiler help}. Note that \cod{-Wall -Wextra} does
\textbf{not} catch the uninitialised \cod{word\_count}, because its address is taken --- GCC gives
up on tracking it. Some bugs need tools beyond warnings: a sanitiser, a unit test, or simply trying
the edge cases deliberately.
\end{concept}

\section{Executing}

\filetag{lytlnybl/user/programs/shell.c}
\begin{lstlisting}[style=c]
bool execute_command(char** words, size_t word_count) {
  if (word_count != 2) return false;      /* every command takes exactly one argument */

  int fd;

  if      (strcmp(words[0], "ls")    == 0) fs_ops(FS_LS,    words[1]);
  else if (strcmp(words[0], "mkdir") == 0) fs_ops(FS_MKDIR, words[1]);
  else if (strcmp(words[0], "touch") == 0) fs_ops(FS_TOUCH, words[1]);
  else if (strcmp(words[0], "rm")    == 0) fs_ops(FS_RM,    words[1]);
  else if (strcmp(words[0], "run")   == 0) fs_ops(FS_RUN,   words[1]);
  else if (strcmp(words[0], "cat")   == 0) { /* ... */ }
  else if (strcmp(words[0], "write") == 0) { /* ... */ }
  else return false;

  return true;
}
\end{lstlisting}

\begin{warn}{Exactly two words, always}
\cod{if (word\_count != 2) return false;} means every command must have exactly one argument. So
\cod{ls} alone fails; you must type \cod{ls /}. There is no way to implement a zero-argument command
such as \cod{help} or \cod{exit} without changing this line.

Combined with the uninitialised counter above, this is why the shell can behave differently between
runs.
\end{warn}

\begin{recipe}{adding commands to the shell}{ch40-shell-command}
Adds \cod{echo} and \cod{rev} as two more branches in \cod{execute\_command} --- which is all a new
command is. Boot it and type \cod{echo hello} or \cod{rev lytlnybl}.

It also adds \cod{help}, deliberately, to show what the \cod{word\_count != 2} rule costs: the
command exists and works, but only if you type \cod{help} followed by a second word, because
anything that is not exactly two words is rejected before the branches are ever reached.
\end{recipe}

\section{\cod{cat}: reading a file}

\begin{lstlisting}[style=c]
else if (strcmp(words[0], "cat") == 0) {
    fd = open(words[1]);
    if (fd < 0) return false;

    char buffer[128];
    int n;

    while ((n = read(fd, buffer, sizeof(buffer))) > 0) {
      write(1, buffer, n);
    }
    printf("\n");
    close(fd);
}
\end{lstlisting}

This is the canonical Unix read loop, and it is correct in shape: read a chunk, write exactly what
was read, stop at zero.

\begin{danger}{But it triggers the kernel bug from Chapter 34}
\cod{write(1, buffer, n)} with \cod{n} up to 128 --- and the kernel's write handler does
\cod{((char*)buffer)[count] = '\textbackslash 0';}. For a file whose content fills the buffer
exactly, the kernel writes a zero byte to \cod{buffer[128]}, one past the end of a 128-byte local
array on the user's stack.

Two bugs meeting: a kernel that writes out of bounds and a caller that uses the full buffer. Either
one alone would be harmless. This is how real security holes are made.
\end{danger}

\section{\cod{write}: the growing buffer}

\begin{lstlisting}[style=c]
size_t capacity = 128;
size_t length = 0;
char* buffer = malloc(capacity);

while (true) {
    char input[128];
    int n = read(0, input, sizeof(input));

    if (n < 0) { free(buffer); close(fd); return false; }
    if (n == 0) break;

    if (length + n > capacity) {
        while (length + n > capacity) capacity *= 2;
        char *new_buffer = realloc(buffer, capacity);
        if (new_buffer == NULL) { free(buffer); close(fd); return false; }
        buffer = new_buffer;
    }

    memcpy(buffer + length, input, n);
    length += n;

    if (input[n-1] == '\n') break;
}
\end{lstlisting}

This one is genuinely well written, and it is worth saying so. It doubles capacity rather than
growing by a constant (amortised $O(1)$ per byte). It checks \cod{realloc} for failure. It keeps
the old pointer until the new allocation succeeds, so a failure does not leak. It frees and closes
on every error path.

That the same file contains both this and \cod{tokenize\_line} is a good illustration of how code
quality varies within a project depending on how much thought a particular function received.

\section{\cod{run}: starting another program}

\filetag{lytlnybl/kernel/interrupts.c}
\begin{lstlisting}[style=c]
case FS_RUN: {
    process_t* child = load_program(path);

    child->parent_pid = current_process->pid;

    current_process->pid_waiting_for = child->pid;
    current_process->state = PROCESS_BLOCKED;
    context_switch(current_process, get_next_process(), regs);
    break;
}
\end{lstlisting}

And the other half, in \cod{exit}:

\begin{lstlisting}[style=c]
case SYSCALL_EXIT: {
    process_t* parent = find_process_by_pid(current_process->parent_pid);

    if (parent != NULL && parent->pid_waiting_for == current_process->pid) {
        parent->pid_waiting_for = 0;
        parent->state = PROCESS_READY;      /* wake the shell */
    }

    current_process->state = PROCESS_TERMINATED;
    context_switch(current_process, get_next_process(), regs);
    break;
}
\end{lstlisting}

\begin{center}
\begin{tikzpicture}[font=\sffamily\scriptsize,
  n/.style={draw=grayborder,rounded corners=2pt,minimum width=4.6cm,minimum height=0.6cm,align=left,inner xsep=7pt}]
  \node[n,fill=greenlight,draw=green] (a) {shell: \texttt{fs\_ops(FS\_RUN, "/bin/user\_test")}};
  \node[n,fill=redlight,draw=red,below=0.22cm of a] (b) {kernel: \texttt{load\_program} creates the child};
  \node[n,fill=redlight,draw=red,below=0.22cm of b] (c) {kernel: child.parent\_pid = shell.pid};
  \node[n,fill=amberlight,draw=amber,below=0.22cm of c] (d) {kernel: shell $\rightarrow$ BLOCKED, switch away};
  \node[n,fill=purplelight,draw=purple,below=0.22cm of d] (e) {child runs, prints, calls \texttt{exit()}};
  \node[n,fill=redlight,draw=red,below=0.22cm of e] (f) {kernel: finds the parent, sets it READY};
  \node[n,fill=greenlight,draw=green,below=0.22cm of f] (g) {shell resumes, prints the prompt};
  \foreach \x/\y in {a/b,b/c,c/d,d/e,e/f,f/g} \draw[-{Stealth[length=2mm]},gray] (\x)--(\y);
\end{tikzpicture}
\end{center}

\begin{concept}{This is \cod{fork} + \cod{exec} + \cod{wait}, collapsed into one}
Unix separates these into three system calls, which is more flexible --- it lets the shell set up
redirection between forking and exec'ing. lytlnyblOS merges them, which is much simpler and enough
for \cod{run}.

The blocking-and-waking mechanism, though, is exactly the real thing: the parent records which child
it waits for, goes to sleep, and the child's exit path wakes it. That handshake is the heart of
process management in every Unix.
\end{concept}

\begin{tryit}[The full round trip, verified]
On a running system:
\begin{lstlisting}[style=txt]
> run /bin/user_test
ab
i have a number it is: 314
malloc: hi
calloc: abc
realloc: hellooo
>
\end{lstlisting}
The child ran, exercised \cod{printf}, \cod{malloc}, \cod{calloc} and \cod{realloc} in user space,
exited, and the shell got its prompt back. Every mechanism in Parts IX through XII, working
together.
\end{tryit}

\begin{danger}{Unless the file does not exist}
Type a path that does not exist --- a typo is enough --- and the missing NULL check from Chapter 33
fires:
\begin{lstlisting}[style=txt]
> run /bin/usertest
(nothing. the prompt never returns. further typing echoes but does nothing.)
\end{lstlisting}
The shell is blocked waiting for a child that was never created, with a pid read out of memory near
address zero. Nothing will ever wake it. This is reproducible on the current code and is the single
easiest serious bug in the project to fix.
\end{danger}

\begin{summary}
The system is complete: it boots, protects itself, manages memory, runs isolated programs,
stores files and gives you a command line. The last part steps back and looks at it as an engineer
rather than a student.
\end{summary}
